\section{The quadratic endpoint of the original envelope}\label{sec:quadratic}
\begin{theorem}\label{thm:quadratic}
Every function in Theorem~\ref{thm:main} is zero-free in
\[
 \Rea s>\sigma_*:=\frac{1507-2\sqrt{921}}{1653}.
\]
\end{theorem}
\begin{proof}
Put $r_0=\sqrt{921}$ and
\begin{equation}\label{eq:quadraticparameters}
 d_*=\frac{49-r_0}{48},\quad
 \ell_*=\frac{33+8r_0}{1653},\quad
 b_*=-\frac4{29}+\frac{230r_0}{26709}.
\end{equation}
We first derive these parameters. At $y=1/2-x=0$, the numerator $N=-108JE_g$ is
\[
 Q_0(\delta)=(-48b+63\ell+21)\delta^2
 +(98b+120\ell-52)\delta-185b/6-925\ell/4+185/4.
\]
The coefficient of $b$ vanishes when
\[
 288\delta^2-588\delta+185=0.
\]
The root in $[0,3/4]$ is $d_*$. Substitution in \eqref{eq:balanced} gives $R_*(d_*,1/2)=2/3$, equivalently $T=d_*-1/3$. Setting the endpoint to zero yields
\[
 \ell_*=\frac{5-6d_*}{9+18d_*}.
\]
The stationary equation $Q_0'(d_*)=0$ then gives
\[
 b_*=\frac{52-(126d_*+120)\ell_*-42d_*}{98-96d_*},
\]
which reduces to \eqref{eq:quadraticparameters}. Also $\sigma_*=11/12-\ell_*/4$.

The rational isolation
\begin{equation}\label{eq:sqrtisolation}
 \frac{30347}{1000}<r_0<\frac{30348}{1000}
\end{equation}
follows from
\[
 921-(30347/1000)^2=59591/10^6>0,\qquad
 (30348/1000)^2-921=1104/10^6>0.
\]
It implies $1/6<\ell_*<17/100$, $1/10<b_*<13/100$, and $d_*>777/2000$.

For the complete rectangle numerator, set
\begin{align*}
 A_0&=-48b_*+63\ell_*+21,\\
 A_1&=60b_*+24\ell_*+54,&B_1&=38b_*+35\ell_*-42,\\
 C_1&=-85b_*/3-425\ell_*/2+85/2,\\
 A_2&=72b_*+84\ell_*+24,&B_2&=-36b_*-10\ell_*+4.
\end{align*}
Expansion gives the identity
\begin{equation}\label{eq:quadraticpoly}
 N=A_0(\delta-d_*)^2+y(A_1\delta^2+B_1\delta+C_1)
                  +y^2(A_2\delta^2+B_2\delta)+96\ell_*y^3\delta^2.
\end{equation}
The double root in the constant coefficient is a consequence of the two displayed elimination equations, not a numerical approximation.

All sign comparisons are in $\Q(r_0)$. Reduce a field expression to $A+Br_0$ using $r_0^2=921$, and use the appropriate endpoint of \eqref{eq:sqrtisolation}. This gives
\[
 A_0>25,\qquad A_1>65,\qquad A_2>46.
\]
For the closest comparisons, the exact field expressions are
\begin{align}
 4A_1(C_1+6/25)-B_1^2
 &=\frac{17317418985607-570341476744r_0}{6990413025}>1,\label{eq:qguard1}\\
 4A_2(3/125)-B_2^2
 &=\frac{-6767947591748+223064581728r_0}{34952065125}>1/30,\label{eq:qguard2}\\
 Q_1'(3/10)&=-2729/145+117028r_0/133545>7,\notag\\
 Q_1(3/10)+1/200&=2140057/66120-40590578r_0/38060325>0,\notag\\
 Q_1(1/3)&=52541/1653-1578082r_0/1522413>3/10,\notag\\
 (3/10)A_2+B_2&=1956/145-6736r_0/133545>11.\notag
\end{align}
Here $Q_1=A_1\delta^2+B_1\delta+C_1$, $Q_2=A_2\delta^2+B_2\delta$, and $Q_3=96\ell_*\delta^2$. Equations~\eqref{eq:qguard1}--\eqref{eq:qguard2} give $Q_1\ge-6/25$, $Q_2\ge-3/125$ by completing squares.

If $0\le\delta\le3/10$, then
\[
 Q_0>25(777/2000-3/10)^2>19/100,\qquad
 N>19/100-3/25-3/500=8/125>0.
\]
If $3/10\le\delta\le1/3$, then
\[
 Q_0>25(777/2000-1/3)^2>3/40.
\]
The derivative guard makes $Q_1$ increasing there, so $Q_1>-1/200$, and the $Q_2$ guard makes $Q_2\ge0$. Thus $N>3/40-1/400=29/400>0$.
If $1/3\le\delta\le3/4$, then $Q_1>0$ and $Q_2,Q_3\ge0$. Hence $N\ge0$, with equality only at $y=0,\delta=d_*$. Since $J>0$, this proves $E_g\le0$ on the entire closed rectangle. Theorem~\ref{thm:geometry} includes the contact point and proves the claimed half-plane.

For comparison,
\[
 \frac78-\sigma_*=\frac{16r_0-485}{13224}>0,\qquad
 \ell_*-\frac{1001}{6000}>\frac1{3306000}>0.
\]
The second strict bound follows by evaluating at the lower endpoint in \eqref{eq:sqrtisolation} and using the strict root inequality. Thus this stage improves both rational stages.
\end{proof}
